% chapters/bagian5-metodologi-penelitian.tex
% proposal.md bagian 5: Metodologi penelitian.
% Contoh teks ilmiah di bawah ini adalah placeholder untuk memperlihatkan
% hasil format; ganti dengan isi proposal yang sebenarnya.

\chapter{Metodologi Penelitian}

\section{Jenis dan Pendekatan Penelitian}
Penelitian ini menggunakan pendekatan eksperimental kuantitatif, yaitu
dengan membangun model klasifikasi sentimen berbasis LSTM,
melatihnya menggunakan data ulasan produk berbahasa Indonesia, dan
mengukur kinerjanya menggunakan metrik evaluasi standar pada bidang
pembelajaran mesin.

\section{Data Penelitian}
Data yang digunakan dalam penelitian ini adalah kumpulan ulasan
produk berbahasa Indonesia beserta label sentimen (positif, negatif,
netral) yang diberikan oleh pengguna melalui penilaian bintang. Data
tersebut direncanakan dibagi menjadi data latih, data validasi, dan
data uji dengan proporsi 70:15:15.

\section{Tahapan Penelitian}
Penelitian ini direncanakan dilaksanakan melalui tahapan berikut.
\begin{enumerate}
  \item \textbf{Studi literatur.} Mengkaji penelitian terdahulu terkait
        klasifikasi sentimen dan metode LSTM \citep{example2021}.
  \item \textbf{Pengumpulan data.} Mengumpulkan data ulasan produk
        beserta label sentimennya.
  \item \textbf{Praproses data.} Melakukan \textit{case folding},
        tokenisasi, normalisasi kata tidak baku, penghapusan stopword,
        dan \textit{stemming}.
  \item \textbf{Perancangan dan pelatihan model.} Merancang arsitektur
        LSTM dan melatihnya menggunakan data latih.
  \item \textbf{Evaluasi model.} Menguji kinerja model menggunakan
        data uji dan membandingkannya dengan model \textit{baseline}.
\end{enumerate}

\section{Jadwal Penelitian}
Penelitian ini direncanakan berlangsung selama satu semester,
mencakup tahapan studi literatur, pengumpulan dan praproses data,
perancangan serta pelatihan model, evaluasi, dan penyusunan laporan
akhir.
